% !TeX root = main.tex
The special unitary group is defined as the group of all unitary matrices $ U \in \mathbb{C}^{N \times N}$ with unit determinant.
First, consider the matrix representation\footnote{A classification of all irreducible representation characterised by the tensor method can be found in \refcite{JakobLinder10.2022}.} of \SU{N}, i.e.\ $ U \in \SU{N} $ being a linear transformation of a vector. Then, the vector $ \psi_{a} = (\psi_{1}, \cdots, \psi_{N}) \in \mathbb{C}^{N} $ of the fundamental representation transforms as
\begin{equation}
	\begin{aligned}
		\psi_{a} \rightarrow \psi^{\prime}_{a} = \tensor{U}{_{a}^{b}} \psi_{b}
		\text{.}
	\end{aligned}
\end{equation}
Note, that the index position of the matrix $ \tensor{U}{_{a}^{b}} := (U)_{ab} $ is \cite[s.][Chap. 8.5 and App. B \& C]{Costa.2012} chosen such that if the first index is a superscript and the second a subscript, i.e.\
$ (\tensor{U}{_{a}^{b}})^{*} = \tensor{(U^{*})}{^{a}_{b}} =: \tensor{U}{^{a}_{b}} $ the complex conjugate matrix is taken. A subscript index, thus correspond to the fundamental representation whereas an superscript index corresponds to the antifundamental representation.
Therefore, a vector $ \phi^{a} = (\phi^{1}, \cdots, \phi^{N})^{\trans}~\in~\mathbb{C}^{N} $ transforming \wrt the anti-fundamental representation transforms as
\begin{equation}
	\begin{aligned}
		\phi^{a} \rightarrow \phi^{\prime a}
		= \tensor{U}{^{a}_{b}} \phi^{b}
		= \phi^{b} \tensor{(U^{\dagger})}{_{b}^{a}}
		\text{.}
	\end{aligned}
\end{equation}
This means, for example, that the hermitian conjugate matrix is just $ \tensor{(U^{\dagger})}{_{a}^{b}} = \tensor{U}{^{b}_{a}} $ and the unitarity condition becomes
\begin{equation}
	\begin{aligned}
		\tensor*{\delta}{_{a}^{d}}
		\equiv \tensor{U}{_{a}^{b}} \tensor{(U^{\dagger})}{_{b}^{d}}
		= \tensor{U}{_{a}^{b}} \tensor{U}{^{d}_{b}}
		= \tensor{U}{_{a}^{b}} \tensor{U}{^{d}_{c}} \tensor*{\delta}{_{b}^{c}}
		\text{,}
	\end{aligned}
\end{equation}
showing also the invariance of the Kronecker delta \wrt the \SU{N} transformation.\par
%
The explicit transformation matrix in the matrix representation can be chosen as follows:
\begin{equation}
	\begin{aligned}
		U = \exp{\mi \eps^{A} T^{A}}
		\text{,}
	\end{aligned}
\end{equation}
where $ \eps^{A} \in \mathbb{R} $ and $ T^{A} $ hermitian. The exponent, $ \mi \eps^{A} T^{A} $, is part of a representation of the Lie algebra of the group. The representation of the Lie algebra is a set of squared hermitian matrices $ T^{A} $, called generators satisfying the commutator relation
\begin{equation}
	\begin{aligned}
		\com{T^{A}}{T^{B}}
		=
		\mi f^{A B C} T^{C}
		\text{.}
	\end{aligned}
\end{equation}
They are normalised as
\begin{equation}
	\begin{aligned}
		\Tr{T^{A} T^{B}} \equiv \tensor{(T^{A})}{_{i}^{j}} \tensor{(T^{B})}{_{j}^{i}} = T_{F} \delta^{A B}
		\quad
		\text{with}
		\quad
		T_{F} = \frac{1}{2}
		\text{.}
	\end{aligned}
\end{equation}
The generators for \SU{2} and \SU{3} are $ \tensor{(T^{I})}{_{i} ^{k}} := \frac{\tensor{(\tau^{I})}{_{i} ^{k}}}{2} $ and $ \tensor{(T^{A})}{_{a} ^{b}} := \frac{\tensor{(\lambda^{A})}{_{a} ^{b}}}{2} $ and are characterised by the Pauli and the Gell-Mann matrices. For clarity, all fundamental indices are written as subscripts and all antifundamental indices as superscripts. The indices of the adjoint representation can be identified by capital letters.\par
The so-called structure constants, $ f^{A B C} $, which are antisymmetric in all indices are
\begin{equation}
	\begin{aligned}
		f^{A B C} = - \frac{\mi}{T_{F}} \Tr{\left[T^{A}, T^{B}\right] T^{C}}
		\text{.}
	\end{aligned} \label{eq:structureconstSUN}
\end{equation}
%~\cite[s.][S.~482--488]{Schwartz2014}
%\sum_{A} \left(T^{A} T^{A}\right)_{i j} = C_{F} \delta_{i j}
%f^{A C D} f^{B C D} = C_{A} \delta^{A B}
Furthermore, the Fierz identity for generators in the fundamental representation
\begin{equation}
	\begin{aligned}
		\sum_{A} \tensor{(T^{A})}{_{i} ^{j}} \tensor{(T^{A})}{_{k} ^{m}}
		= T_{F} \big(
		\tensor*{\delta}{_{i} ^{m}} \tensor*{\delta}{_{k} ^{j}}
		- \frac{1}{N} \tensor*{\delta}{_{i} ^{j}} \tensor*{\delta}{_{k} ^{m}}
		\big)
		\text{,}
	\end{aligned}
\end{equation}
is satisfied. Note, that the generator of the fundamental representation $ \tensor{(T^{A})}{_{i} ^{j}} $ transforms with one fundamental and one antifundamental index.\par
%
\subsection{\eps{} - tensors}
The epsilon tensor of \SU{N} with fundamental indices is defined by
\begin{equation}
	\begin{aligned}
		\tensor{\eps}{_{i_{1} \ldots i_{N}}} \equiv
		\det{
			\left(
			\begin{array}{ccc}
				\tensor*{\delta}{_{i_{1} 1}} & \cdots & \tensor*{\delta}{_{i_{1} N}} \\
				\vdots & & \vdots \\
				\tensor*{\delta}{_{i_{N} 1}} & \cdots & \tensor*{\delta}{_{i_{N} N}}
			\end{array}
			\right)
		}
		\text{.}
	\end{aligned}
\end{equation}
The definition for the epsilon tensor with anti-fundamental indices is similar. Furthermore, two epsilon tensors can be expressed by Kronecker deltas:
\begin{equation}
	\begin{aligned}
		\tensor{\eps}{_{i_{1} \ldots i_{N}}} \tensor{\eps}{^{j_{1} \ldots j_{N}}}
		=
		\det{
			\left(
			\begin{array}{ccc}
				\tensor*{\delta}{_{i_{1} j_{1}}} & \cdots & \tensor*{\delta}{_{i_{1} j_{N}}} \\
				\vdots & & \vdots \\
				\tensor*{\delta}{_{i_{N} j_{1}}} & \cdots & \tensor*{\delta}{_{i_{N} j_{N}}}
			\end{array}
			\right)
		}
		\text{,}
	\end{aligned}
\end{equation}
yielding for \SU{2}
\begin{equation}
	\begin{aligned}
		\tensor{\eps}{_{i j}} \tensor{\eps}{^{k m}}
		= \tensor*{\delta}{_{i}^{k}} \tensor*{\delta}{_{j}^{m}}
		- \tensor*{\delta}{_{i}^{m}} \tensor*{\delta}{_{j}^{k}}
		\text{,} \quad
		\tensor{\eps}{_{i j}} \tensor{\eps}{^{i k}}
		= \tensor*{\delta}{_{j}^{k}}
		\text{,} \quad
		\tensor{\eps}{_{i j}} \tensor{\eps}{^{i j}}
		= 2
	\end{aligned} \label{eq:epsilonSU2}
\end{equation}
and \SU{3}:
\begin{equation}
	\begin{aligned}
		\tensor{\eps}{_{i j k}} \tensor{\eps}{^{i m n}}
		= \tensor*{\delta}{_{j}^{m}} \tensor*{\delta}{_{k}^{n}}
		- \tensor*{\delta}{_{j}^{n}} \tensor*{\delta}{_{k}^{m}}
		\text{,} \quad
		\tensor{\eps}{_{i j k}} \tensor{\eps}{^{i j m}}
		= 2 \tensor*{\delta}{_{k}^{m}}
		\text{,} \quad
		\tensor{\eps}{_{i j k}} \tensor{\eps}{^{i j k}}
		= 6
		\text{.}
	\end{aligned} \label{eq:epsilonSU3}
\end{equation}
Note, especially that the signs for \SU{2} in \cref{eq:epsilonSU2} differ from those in the \SLtwoC{} notation in \cref{eq:2epsilonsSL2C}. This is the case, because $ \tensor{\eps}{_{\alpha \beta}} $ is defined as the inverse of $ \tensor{\eps}{^{\alpha \beta}} $ in \cref{eq:SL2Cepsilon}, whereas $ \tensor{\eps}{_{i j}} $ and $ \tensor{\eps}{^{i j}} $ for \SU{2} are numerically identical.
%
Since all irreducible representations of \SU{N} and thus all invariant operators can easily be found by considering only fundamental indices, all fields\footnote{The transformation properties of each field are listed in \cref{tab:fields}.} should be expressed in the fundamental representation.
The fundamental indices can then be identified by a \gls{yt} indicating the corresponding irreducible representation of the field. \par
%
\subsection{Index conversion}
Antifundamental indices are converted with the $ \epsilon $-tensor of the corresponding group:
\begin{equation}
	\boxed{
		\begin{aligned}
			\eps_{i j} \tensor{(\psi^{\dagger})}{^{j}}
			=& \tensor{(\psi^{\dagger})}{_{i}}
			&\sim \ytableaushort{i} \text{, for }
			\psi^{\dagger} \in \{H^{\dagger}, L^{\dagger}, Q^{\dagger}\}\\
			%		
			\eps_{a b c} \tensor{(\psi^{\dagger})}{^c}
			=& \psi^{\dagger}_{a b}
			&\sim \ytableaushort{a,b}
			 \text{, for }
			\psi^{\dagger} \in \{Q^{\dagger}, u_{\mathbb{C}}, d_{\mathbb{C}}\} \text{.}
		\end{aligned} \label{eq:afundSU2andSU3}
	}
\end{equation}
%
A field $ F^{A} $ of the adjoint representation of \SU{N} is contracted with the corresponding generator $ \tensor{(T^{A})}{_{b} ^{a}} $, and the antifundamental index is again converted with an $ \epsilon $-tensor \cite{Li2021c}:
% Sec 3.1
\begin{equation}
	\ytableausetup{boxsize=2em}
	\begin{aligned}
		\eps_{a_{1} \ldots a_{N}} \tensor{(T^{A})}{_{b} ^{a_{N}}} F^{A} =: F_{a_{1} \ldots a_{N-1} b}
		\sim
		\begin{ytableau}
			a_{1} & b \\
			\none[\vdots] \\
			\scriptstyle a_{N-1}
		\end{ytableau}
	\end{aligned}
	\ytableausetup{boxsize=normal, smalltableaux}
\end{equation}
{\color{red}
In order to follow the conventions of \refcite{Li2021}, an additional factor $ 2 $ is multiplied for the \SU{2} and \SU{3} field strength tensors:
\begin{equation}
	\boxed{
		\begin{aligned}
			2 \eps_{j k} \tensor{(T^{I})}{_{i} ^{k}} W^{I}
			&= W_{(i j)}
			&\sim \ytableaushort{ij} \text{,}\\
			2 \eps_{a c d} \tensor{(T^{A})}{_{b} ^{d}} G^{A}
			&= G_{a b c}
			&\sim \ytableaushort{ab,c}
			\text{.}
		\end{aligned} \label{eq:adjointSU2andSU3}
	}
\end{equation}
}
Note that the indices of both field strength tensors obey so-called internal symmetries, which can be read from the \glspl{yt}. The $ W $, for example, is symmetric in its fundamental \SU{2} indices.
%

\paragraph{Summary}
The \gls{sm} fields containing different than fundamental gauge indices are rewritten with only fundamental ones like shown in the following:
\begin{equation}
	\boxed{
		\begin{alignedat}{5}
			\eps_{i j} \tensor{(H^{\dagger})}{^{j}} &= \tensor{(H^{\dagger})}{_{i}}
			&&\sim \ytableaushort{i}
			&&\Leftrightarrow
			&&\tensor{(H^{\dagger})}{^{i}} &&= \tensor{(H^{\dagger})}{_{j}} \eps^{j i}
			\text{,} \\
			%			
			2 \eps_{j k} \tensor{(T^{I})}{_{i} ^{k}} W^{I} &= W_{(i j)}
			&&\sim \ytableaushort{ij}
			&&\Leftrightarrow
			&& W^{I} && = \tensor{(T^{I})}{_{k} ^{i}} W_{(i j)} \eps^{j k}
			\text{,} \\
			%
			\eps_{a b c} \tensor{(Q^{\dagger})}{^c} &= Q_{a b}^{\dagger}
			&&\sim \ytableaushort{a,b}
			&&\Leftrightarrow
			&& \tensor{(Q^{\dagger})}{^a} &&= \frac{1}{2} \eps^{a b c} Q^{\dagger}_{b c}
			\text{,} \\
			%
			2 \eps_{a c d} \tensor{(T^{A})}{_{b} ^{d}} G^{A} &= G_{a b c}
			&&\sim \ytableaushort{ab,c}
			&&\Leftrightarrow
			&& G^{A} &&= \frac{\tensor{(T^{A})}{_{d} ^{b}}}{2} \eps^{a c d} G_{a b c}
			\text{,}
		\end{alignedat}
	}
\end{equation}
where $ W_{(i k)} := (W_{i k} + W_{k i})/ 2 $ is the symmetric part of $ W $.
\todo{Show \cref{eq:adjointSU2andSU3} for an adjoint tensor of every \SU{N}.}